\documentclass[journal]{IEEEtran}

\usepackage{amsmath,amsfonts,amssymb}
\usepackage{graphicx}
\usepackage{cite}
\usepackage{booktabs}
\usepackage{multirow}
\usepackage{url}
\usepackage{hyperref}
\usepackage{xcolor}

\hypersetup{
    colorlinks=true,
    linkcolor=blue!70!black,
    citecolor=blue!70!black,
    urlcolor=blue!70!black
}

\begin{document}

\title{Accuracy, Faithfulness, and Robustness: An Explainable Benchmark of CNNs, Vision Transformers, and Vision-Language Models for Scene Recognition}

\author{MD~Matiur~Rahman%
% TODO: co-author permission (Prof. Huanjie Tao)
\thanks{MD Matiur Rahman is with the School of Information Science and Technology, Nantong University, Nantong 226019, China (e-mail: matiurrahman1353@gmail.com).}%
\thanks{Source code, trained checkpoints, and reproducibility configurations are publicly available at: \url{https://github.com/matiur-rahman1/explainable-scene-recognition}.}}

\markboth{IEEE Access}{Rahman: Accuracy, Faithfulness, and Robustness in Explainable Scene Recognition}

\maketitle

\begin{abstract}
Multi-class scene recognition is a foundational task in computer vision supporting autonomous navigation, robotic perception, and environmental monitoring. Although modern Convolutional Neural Networks (CNNs), hierarchical Vision Transformers (ViTs), and contrastive Vision-Language Models (VLMs) achieve high classification accuracy on natural scene benchmarks, their deployment in safety-critical applications is hindered by two persistent challenges: opaque decision-making and vulnerability to adversarial perturbations. Existing studies typically evaluate predictive accuracy, visual explainability, or adversarial robustness in isolation, leaving their cross-architectural trade-offs poorly understood. In this paper, we present a unified, reproducible benchmark comparing representative CNN backbones (ResNet-50, EfficientNet-B0), a hierarchical Vision Transformer (Swin-Tiny), and a contrastive vision-language foundation model (CLIP ViT-B/16 in zero-shot and linear-probe settings) across natural scene benchmarks (Intel Natural Scenes and a 100-class SUN397 protocol). Beyond standard classification accuracy, macro-F1, and inference latency, we quantitatively evaluate three gradient-based class activation mapping techniques---Grad-CAM, Grad-CAM++, and LayerCAM---using progressive Insertion and Deletion Area Under the Curve (AUC) faithfulness metrics, and assess adversarial resilience under Fast Gradient Sign Method (FGSM) and Projected Gradient Descent ($\text{PGD-}L_\infty$) attacks as well as PGD adversarial training (PGD-AT). Our empirical results demonstrate a pronounced trilemma: while Swin-Tiny and linear-probe CLIP achieve superior classification accuracy (94.60\% and \textbf{[TBD]}\% on Intel Natural Scenes), ResNet-50 yields significantly higher spatial explanation faithfulness under Grad-CAM (0.8405 Insertion AUC) and stronger native PGD resilience (19.17\% vs.\ 1.67\% at $\epsilon=0.03$). Furthermore, second-order gradient weighting in Grad-CAM++ substantially improves explanation faithfulness on Swin-Tiny compared to standard Grad-CAM (0.7666 vs.\ 0.6031 Insertion AUC), and PGD-AT restores adversarial accuracy while regularizing attention drift.
\end{abstract}

\begin{IEEEkeywords}
Scene recognition, explainable artificial intelligence (XAI), Grad-CAM, Vision Transformer, Swin Transformer, CLIP, insertion and deletion AUC, adversarial robustness, PGD adversarial training.
\end{IEEEkeywords}

\section{Introduction}
\IEEEPARstart{V}{isual} scene recognition---the task of categorizing an image according to its holistic environmental context (e.g., \emph{glacier}, \emph{forest}, \emph{street}, or \emph{indoor abbey})---is a cornerstone of scene understanding in computer vision \cite{oliva2001modeling,lazebnik2006beyond,xiao2010sun,zhou2017places}. Unlike object-centric classification \cite{deng2009imagenet}, where a single foreground entity typically dominates the decision boundary, natural scene recognition requires integrating distributed spatial cues, background textures, and structural layouts across multiple scales \cite{quattoni2009recognizing,xiao2016sun}. Deep neural networks have driven rapid progress on scene benchmarks, evolving from deep residual and compound-scaled convolutional networks \cite{he2016deep,tan2019efficientnet,liu2022convnet} to self-attention Vision Transformers \cite{dosovitskiy2021an,touvron2021training,liu2021swin} and web-scale contrastive Vision-Language Models (VLMs) such as CLIP \cite{radford2021learning,jia2021scaling,cherti2023reproducible}.

Despite their impressive top-1 accuracy, deploying deep scene classifiers in high-stakes environments such as autonomous driving, unmanned aerial surveillance, and robotic navigation raises two fundamental concerns:
\begin{enumerate}
    \item \textbf{Lack of Quantitative Interpretability:} Deep vision backbones operate as high-dimensional black boxes. While post-hoc attribution algorithms such as Grad-CAM \cite{selvaraju2017grad}, Grad-CAM++ \cite{chattopadhay2018grad}, and LayerCAM \cite{jiang2021layercam} produce plausible saliency overlays, qualitative heatmaps alone can be misleading \cite{adebayo2018sanity,ghorbani2019interpretation}. In our prior collaborative work on explainable machine learning for phishing detection \cite{tulla2025phishing}, we demonstrated on a balanced 10,000-URL dataset with 48 lexical and content features that pairing predictive classifiers with rigorous post-hoc explainability is essential for bridging technical efficacy and operational trustworthiness. Extending this principle to high-dimensional natural images requires validating visual explanations \emph{quantitatively} via pixel-perturbation faithfulness metrics \cite{samek2016evaluating,petsiuk2018rise}.
    \item \textbf{Fragility Under Adversarial Perturbations:} Deep classifiers are notoriously vulnerable to imperceptible $L_\infty$-bounded perturbations generated by gradient-based attacks such as FGSM \cite{goodfellow2015explaining} and PGD \cite{madry2018towards}. Moreover, adversarial noise does not merely flip predicted labels; it induces severe \emph{attention drift} in class activation maps, shifting model focus away from semantically meaningful regions toward background artifacts \cite{ghorbani2019interpretation,tsipras2019robustness}.
\end{enumerate}

Crucially, existing literature largely investigates accuracy, explainability faithfulness, and adversarial robustness along separate tracks. Comparative studies of CNNs and Vision Transformers \cite{bhojanapalli2021understanding,bai2021transformers} examine clean and adversarial accuracy without quantifying whether their internal spatial attributions remain faithful. Conversely, XAI evaluations \cite{selvaraju2017grad,chattopadhay2018grad,jiang2021layercam,chefer2021transformer} rarely benchmark CNNs, hierarchical Swin Transformers, and multimodal CLIP encoders under a unified pixel-perturbation and adversarial training protocol.

To bridge this gap, we conduct a systematic, multi-seed experimental investigation across convolutional, hierarchical transformer, and vision-language paradigms. The primary contributions of this work are fourfold:
\begin{itemize}
    \item \textbf{Unified Cross-Paradigm Scene Benchmark:} We benchmark CNNs (ResNet-50 \cite{he2016deep}, EfficientNet-B0 \cite{tan2019efficientnet}), a shifted-window Vision Transformer (Swin-Tiny \cite{liu2021swin}), and a contrastive vision-language backbone (CLIP ViT-B/16 \cite{radford2021learning} in both prompt-ensembled zero-shot and linear-probe configurations) across multiple random seeds on Intel Natural Scenes and a deterministic 100-class SUN397 protocol \cite{xiao2010sun}, accompanied by paired McNemar statistical significance tests \cite{mcnemar1947note}.
    \item \textbf{Quantitative XAI Faithfulness Analysis:} We evaluate Grad-CAM \cite{selvaraju2017grad}, Grad-CAM++ \cite{chattopadhay2018grad}, and LayerCAM \cite{jiang2021layercam} across CNN and Transformer token grids using progressive Insertion and Deletion AUC metrics \cite{petsiuk2018rise}, uncovering why second-order gradient weighting (Grad-CAM++) is distinctly advantageous for shifted-window self-attention representations.
    \item \textbf{Adversarial Robustness, Label-Leaking Diagnostics, and PGD-AT:} We evaluate all architectures under untargeted FGSM, predicted-label FGSM, and multi-step $\text{PGD-}L_\infty$ attacks \cite{goodfellow2015explaining,kurakin2017adversarial,madry2018towards}, diagnose the single-step FGSM label-leaking phenomenon induced by label smoothing, and quantify how PGD adversarial training (PGD-AT) impacts both classification resilience and CAM spatial faithfulness.
    \item \textbf{Reproducible Open-Source Framework:} We release a modular PyTorch framework, automated LaTeX table generators, checkpoint resume pipelines, and an interactive Gradio diagnostic interface to facilitate reproducible research.
\end{itemize}

The remainder of this paper is organized as follows. Section~\ref{sec:related_work} reviews related work in scene recognition, vision backbones, XAI faithfulness, and adversarial robustness. Section~\ref{sec:methodology} formalizes the architectures, gradient-based attribution methods, faithfulness metrics, and adversarial protocols. Section~\ref{sec:setup} details the datasets and experimental setup. Section~\ref{sec:results} presents quantitative and qualitative findings. Section~\ref{sec:discussion} discusses the accuracy--faithfulness--robustness trade-off and practical limitations, and Section~\ref{sec:conclusion} concludes with future directions including open-vocabulary grounded detection.

\section{Related Work}
\label{sec:related_work}

\subsection{Natural Scene Recognition Benchmarks}
Early scene recognition research progressed from holistic spatial envelope descriptors (GIST) \cite{oliva2001modeling} and spatial pyramid matching over local features \cite{lazebnik2006beyond} to indoor scene representations \cite{quattoni2009recognizing}. The introduction of the Scene UNderstanding (SUN397) database by Xiao et al.\ \cite{xiao2010sun,xiao2016sun} established a comprehensive 397-category benchmark spanning indoor, urban, and natural environments. Subsequently, Zhou et al.\ \cite{zhou2017places} introduced Places365, scaling scene data to 10 million images and demonstrating that scene-centric convolutional features differ structurally from object-centric ImageNet representations \cite{deng2009imagenet,zhou2016learning}. For rapid architectural ablation and clean natural-scene boundary analysis, the 6-class Intel Natural Scenes benchmark (\emph{buildings}, \emph{forest}, \emph{glacier}, \emph{mountain}, \emph{sea}, \emph{street}) provides high intra-class diversity and natural semantic ambiguities (e.g., \emph{glacier} vs.\ \emph{mountain}).

\subsection{CNNs, Vision Transformers, and Vision-Language Models}
Convolutional architectures rely on local receptive fields, translation equivariance, and hierarchical spatial pooling. ResNet-50 \cite{he2016deep} introduced residual skip connections enabling deep feature extraction, while EfficientNet-B0 \cite{tan2019efficientnet} optimized depth, width, and resolution via compound scaling and inverted residual blocks. ConvNeXt \cite{liu2022convnet} later modernized pure CNN design principles. In contrast, the Vision Transformer (ViT) \cite{dosovitskiy2021an} and DeiT \cite{touvron2021training} process images as sequences of non-overlapping $16\times16$ patches via global multi-head self-attention. To restore hierarchical multi-scale representations and linear computational complexity with respect to image size, Liu et al.\ \cite{liu2021swin} proposed the Swin Transformer, computing self-attention within local windows and shifting window partitions across consecutive layers.

More recently, contrastive Vision-Language Models such as CLIP \cite{radford2021learning}, ALIGN \cite{jia2021scaling}, and OpenCLIP \cite{cherti2023reproducible} align image and text embeddings across hundreds of millions of image--caption pairs. By encoding class names into natural language prompt templates (e.g., \emph{``a photo of a \{scene\}''}) \cite{radford2021learning,zhou2022learning}, CLIP enables zero-shot transfer to downstream scene categories without task-specific gradient updates, as well as data-efficient linear probing on frozen visual embeddings.

\subsection{Gradient-Based Visual Explainability and Faithfulness}
Post-hoc explainability methods aim to attribute model predictions to input features or spatial regions \cite{ribeiro2016should,lundberg2017unified,sundararajan2017axiomatic,tulla2025phishing}. In prior work by the authors \cite{tulla2025phishing}, post-hoc feature attribution was integrated with Decision Tree, XGBoost, and Logistic Regression classifiers to establish interpretable and legally accountable phishing detection. In computer vision, Class Activation Mapping (CAM) \cite{zhou2016learning} and its gradient-weighted generalization Grad-CAM \cite{selvaraju2017grad} compute channel-wise importance weights by global average pooling the gradients of the target class score with respect to the final spatial feature map. Because global average pooling can underweight small or multiple object instances, Chattopadhay et al.\ \cite{chattopadhay2018grad} formulated Grad-CAM++, incorporating second- and third-order partial derivatives to weight pixel-wise positive gradients. LayerCAM \cite{jiang2021layercam} further refines spatial resolution by weighting activations element-wise via positive gradients across intermediate or final layers, avoiding coarse pooling artifacts, while Score-CAM \cite{wang2020score} eliminates gradients at the cost of hundreds of forward passes. For Vision Transformers, attribution requires either reshaping patch token activations back into 2D spatial grids or aggregating attention rollout and relevance propagation across layers \cite{abnar2020quantifying,chefer2021transformer}.

Crucially, Adebayo et al.\ \cite{adebayo2018sanity} showed that visual inspection of saliency maps is insufficient because certain attribution methods act as parameter-insensitive edge detectors. Consequently, Samek et al.\ \cite{samek2016evaluating} and Petsiuk et al.\ \cite{petsiuk2018rise} established quantitative faithfulness metrics---specifically progressive pixel \emph{Deletion} and \emph{Insertion} curves---which measure how rapidly target class probability collapses when the highest-attributed pixels are removed, or recovers when they are revealed on a uninformative baseline image.

\subsection{Adversarial Robustness in Vision Backbones}
Szegedy et al.\ \cite{szegedy2014intriguing} and Goodfellow et al.\ \cite{goodfellow2015explaining} demonstrated that small, worst-case $L_\infty$ perturbations generated via the Fast Gradient Sign Method (FGSM) severely degrade neural network accuracy. Kurakin et al.\ \cite{kurakin2017adversarial} scaled adversarial evaluation to ImageNet and identified the \emph{label leaking} phenomenon, wherein single-step FGSM perturbations constructed using the ground-truth label inadvertently regularize or overshoot the decision boundary at larger step sizes $\epsilon$. Madry et al.\ \cite{madry2018towards} formalized multi-step Projected Gradient Descent (PGD) as a universal first-order adversary and established PGD adversarial training (PGD-AT) via min-max optimization. Subsequent theoretical and empirical work by Tsipras et al.\ \cite{tsipras2019robustness} and Zhang et al.\ \cite{zhang2019theoretically} characterized the fundamental trade-off between standard clean accuracy and adversarial robustness, while Bhojanapalli et al.\ \cite{bhojanapalli2021understanding} and Bai et al.\ \cite{bai2021transformers} compared the robustness profiles of ViTs and CNNs. Furthermore, Ghorbani et al.\ \cite{ghorbani2019interpretation} showed that adversarial perturbations destabilize interpretation maps, motivating joint analysis of robustness and XAI faithfulness.

\section{Methodology}
\label{sec:methodology}

\subsection{Problem Formulation and Architectures}
Let $\mathcal{D} = \{(x_i, y_i)\}_{i=1}^{N}$ denote a multi-class natural scene dataset, where $x_i \in [0, 1]^{3 \times H \times W}$ is an RGB image resized to $H=W=224$ and $y_i \in \{1, \dots, K\}$ is its scene category label. An input image is standardized channel-wise using ImageNet statistics $\mu = (0.485, 0.456, 0.406)$ and $\sigma = (0.229, 0.224, 0.225)$ via $\tilde{x} = (x - \mu) / \sigma$. A deep classifier $f_\theta: \mathbb{R}^{3 \times H \times W} \to \mathbb{R}^K$ maps $\tilde{x}$ to class logits $z = f_\theta(\tilde{x})$, yielding posterior probabilities $p_c(\tilde{x}) = \text{softmax}(z)_c = \exp(z_c) / \sum_{j=1}^K \exp(z_j)$.

We benchmark five model configurations spanning three architectural paradigms:
\begin{enumerate}
    \item \textbf{ResNet-50 \cite{he2016deep}:} A 50-layer residual CNN comprising four bottleneck stages with $3\times3$ and $1\times1$ convolutions, terminating in a $7\times7\times2048$ activation tensor at \texttt{layer4[-1]} prior to global average pooling.
    \item \textbf{EfficientNet-B0 \cite{tan2019efficientnet}:} A mobile-sized compound-scaled CNN utilizing mobile inverted bottleneck (MBConv) blocks and squeeze-and-excitation attention, producing a $7\times7\times1280$ feature tensor at \texttt{conv\_head}.
    \item \textbf{Swin-Tiny \cite{liu2021swin}:} A four-stage hierarchical Vision Transformer that partitions images into $4\times4$ patches and applies Window Multi-Head Self-Attention (W-MSA) and Shifted-Window MSA (SW-MSA) within $7\times7$ local windows. Its final stage outputs a token grid of spatial resolution $7\times7$ with channel dimension $C=768$.
    \item \textbf{CLIP ViT-B/16 (Zero-Shot) \cite{radford2021learning}:} A contrastive vision-language model with a 12-layer Vision Transformer image encoder $g_I(\cdot)$ ($16\times16$ patches, $14\times14 = 196$ spatial tokens plus one \texttt{[CLS]} token, $C=768$) and a 12-layer Transformer text encoder $g_T(\cdot)$. For each scene category $c \in \{1,\dots,K\}$, we construct an ensemble of $M=5$ prompt templates $\{t_{c,m}\}_{m=1}^M$ (e.g., \emph{``a photo of a \{class\}.''}, \emph{``a photo of the \{class\} scene.''}) and compute the L2-normalized prompt-ensembled text prototype:
    \begin{equation}
        \bar{w}_c = \frac{\sum_{m=1}^{M} \hat{g}_T(t_{c,m})}{\left\| \sum_{m=1}^{M} \hat{g}_T(t_{c,m}) \right\|_2}, \quad \hat{g}_T(t) = \frac{g_T(t)}{\|g_T(t)\|_2}.
    \end{equation}
    Given an image embedding $\hat{v} = g_I(\tilde{x}_{\text{clip}}) / \|g_I(\tilde{x}_{\text{clip}})\|_2$, where $\tilde{x}_{\text{clip}}$ is differentiably re-normalized to CLIP statistics, the zero-shot logit for class $c$ is:
    \begin{equation}
        z_c^{\text{ZS}}(\tilde{x}) = \tau \cdot (\hat{v}^\top \bar{w}_c),
    \end{equation}
    where $\tau = \exp(\text{logit\_scale})$ is the learned temperature parameter.
    \item \textbf{CLIP ViT-B/16 (Linear Probe):} Freezes the pretrained CLIP visual encoder $g_I(\cdot)$ and trains a linear classification head $W \in \mathbb{R}^{K \times D}, b \in \mathbb{R}^K$ on top of the L2-normalized visual representation $\hat{v}$, warm-started from the zero-shot text prototypes $W_0 = \tau [\bar{w}_1, \dots, \bar{w}_K]^\top$.
\end{enumerate}

Supervised fine-tuning (for ResNet-50, EfficientNet-B0, Swin-Tiny, and Linear-Probe CLIP) minimizes the label-smoothed cross-entropy loss ($\varepsilon_{\text{ls}} = 0.05$):
\begin{equation}
    \mathcal{L}_{\text{CE}}(f_\theta(\tilde{x}), y) = -\sum_{c=1}^{K} \left[ (1 - \varepsilon_{\text{ls}})\mathbb{I}(c = y) + \frac{\varepsilon_{\text{ls}}}{K} \right] \log p_c(\tilde{x}).
    \label{eq:ls_ce}
\end{equation}

\subsection{Gradient-Based Visual Explanations}
Let $A \in \mathbb{R}^{C \times u \times v}$ denote the spatial feature tensor at the target layer of a backbone. For CNNs, $A$ is directly the 2D convolutional activation map ($u=v=7$). For Swin-Tiny, the final normalization layer activation $T \in \mathbb{R}^{(u \cdot v) \times C}$ is reshaped into $A \in \mathbb{R}^{C \times 7 \times 7}$. For CLIP ViT-B/16, the \texttt{[CLS]} token is discarded from the final residual block layer-norm output $T \in \mathbb{R}^{(1 + u \cdot v) \times C}$, and the remaining $196$ patch tokens are reshaped into a 2D spatial tensor $A \in \mathbb{R}^{C \times 14 \times 14}$. Let $y^c = z_c$ be the pre-softmax logit for target class $c$, and let $G_{k,i,j}^c = \frac{\partial y^c}{\partial A_{k,i,j}}$ denote its partial derivative at channel $k$ and spatial location $(i,j)$.

\subsubsection{Grad-CAM}
Grad-CAM \cite{selvaraju2017grad} computes channel importance weights $\alpha_k^c$ via global average pooling of the first-order gradients:
\begin{equation}
    \alpha_k^c = \frac{1}{u \cdot v} \sum_{i=1}^{u} \sum_{j=1}^{v} \frac{\partial y^c}{\partial A_{k,i,j}}, \quad
    L_{\text{Grad-CAM}}^c(i,j) = \text{ReLU}\!\left( \sum_{k=1}^{C} \alpha_k^c A_{k,i,j} \right).
\end{equation}

\subsubsection{Grad-CAM++}
To prevent small or spatially dispersed visual cues from being diluted by uniform spatial averaging, Grad-CAM++ \cite{chattopadhay2018grad} introduces pixel-wise weighting coefficients $\beta_{k,i,j}^c$ derived from second- and third-order gradient terms:
\begin{equation}
    \beta_{k,i,j}^c = \frac{\left(\frac{\partial y^c}{\partial A_{k,i,j}}\right)^2}{2\left(\frac{\partial y^c}{\partial A_{k,i,j}}\right)^2 + \sum_{a=1}^{u}\sum_{b=1}^{v} A_{k,a,b} \left(\frac{\partial y^c}{\partial A_{k,a,b}}\right)^3},
\end{equation}
\begin{equation}
    \alpha_k^{c,\text{++}} = \sum_{i=1}^{u} \sum_{j=1}^{v} \beta_{k,i,j}^c \, \text{ReLU}\!\left(\frac{\partial y^c}{\partial A_{k,i,j}}\right),
\end{equation}
with $L_{\text{Grad-CAM++}}^c(i,j) = \text{ReLU}\!\left(\sum_{k=1}^{C} \alpha_k^{c,\text{++}} A_{k,i,j}\right)$.

\subsubsection{LayerCAM}
LayerCAM \cite{jiang2021layercam} preserves fine-grained spatial localization by multiplying each activation element directly by its thresholded positive gradient prior to channel summation:
\begin{equation}
    L_{\text{LayerCAM}}^c(i,j) = \text{ReLU}\!\left( \sum_{k=1}^{C} \text{ReLU}\!\left(\frac{\partial y^c}{\partial A_{k,i,j}}\right) A_{k,i,j} \right).
\end{equation}
Each resulting map $L^c \in \mathbb{R}^{u \times v}$ is bilinearly upsampled to $(H, W)$ and min-max normalized to $S^c \in [0, 1]^{H \times W}$.

\subsection{Quantitative XAI Faithfulness: Insertion and Deletion AUC}
To evaluate whether the high-attribution pixels in $S^c$ genuinely drive the model's prediction $c = \arg\max_j p_j(\tilde{x})$, we employ progressive pixel perturbation \cite{petsiuk2018rise,samek2016evaluating}. Let $\pi$ be the permutation of pixel coordinates sorted in descending order of attribution $S^c$, and let $\tilde{x}_{\text{blur}} = \text{AvgPool}_{21\times 21}(\tilde{x})$ be a heavily blurred baseline image that removes discriminative high-frequency structures without introducing unnatural black-pixel boundary edges. For $S=20$ uniform perturbation fractions $\gamma_s = s / S \in [0, 1]$ ($s=0,\dots,S$), let $M(\gamma_s) \in \{0, 1\}^{H \times W}$ be the binary mask setting the top $\lfloor \gamma_s H W \rfloor$ pixels in $\pi$ to $1$. The progressive Insertion and Deletion images are defined as:
\begin{align}
    \tilde{x}_{\text{ins}}(\gamma_s) &= M(\gamma_s) \odot \tilde{x} + (1 - M(\gamma_s)) \odot \tilde{x}_{\text{blur}}, \\
    \tilde{x}_{\text{del}}(\gamma_s) &= (1 - M(\gamma_s)) \odot \tilde{x} + M(\gamma_s) \odot \tilde{x}_{\text{blur}}.
\end{align}
The Insertion AUC ($\uparrow$ higher is better) and Deletion AUC ($\downarrow$ lower is better) are computed via the trapezoidal rule:
\begin{equation}
    \text{AUC}_{\text{ins}} = \sum_{s=1}^{S} \frac{p_c(\tilde{x}_{\text{ins}}(\gamma_{s-1})) + p_c(\tilde{x}_{\text{ins}}(\gamma_s))}{2 S},
\end{equation}
and analogously for $\text{AUC}_{\text{del}}$.

\subsection{Adversarial Attacks and PGD Adversarial Training}
We evaluate white-box adversarial robustness under $L_\infty$-norm constraints of radius $\epsilon \in \{0.01, 0.03, 0.05\}$ in raw $[0, 1]$ pixel space. Because the network operates on normalized inputs $\tilde{x}$, the per-channel perturbation budget and valid domain bounds in normalized space are $\tilde{\epsilon} = \epsilon / \sigma$, $\tilde{l} = (0 - \mu)/\sigma$, and $\tilde{u} = (1 - \mu)/\sigma$.

\subsubsection{FGSM and Label-Leaking Diagnostics}
Standard single-step FGSM \cite{goodfellow2015explaining} perturbs $\tilde{x}$ along the sign of the loss gradient with respect to the true label $y$:
\begin{equation}
    \tilde{x}_{\text{FGSM}} = \Pi_{[\tilde{l}, \tilde{u}]}\!\left( \tilde{x} + \tilde{\epsilon} \cdot \text{sign}\!\left( \nabla_{\tilde{x}} \mathcal{L}_{\text{CE}}(f_\theta(\tilde{x}), y) \right) \right).
\end{equation}
To diagnose the label-leaking effect \cite{kurakin2017adversarial} on label-smoothed models, we also evaluate predicted-label FGSM ($\text{FGSM}_{\text{pred}}$), which replaces $y$ with the model's clean prediction $\hat{y} = \arg\max_c f_\theta(\tilde{x})_c$.

\subsubsection{Iterative $\text{PGD-}L_\infty$}
Multi-step PGD \cite{madry2018towards} initializes $\tilde{x}^{(0)} = \Pi_{[\tilde{l}, \tilde{u}]}(\tilde{x} + \delta_0)$ with $\delta_0 \sim \mathcal{U}(-\tilde{\epsilon}, \tilde{\epsilon})$ and iterates for $T=10$ steps with scaled step size $\tilde{\alpha} = 2.5 \tilde{\epsilon} / T$:
\begin{equation}
    \tilde{x}^{(t+1)} = \Pi_{[\tilde{l}, \tilde{u}]}\!\left( \Pi_{\mathcal{B}_\infty(\tilde{x}, \tilde{\epsilon})}\!\left( \tilde{x}^{(t)} + \tilde{\alpha} \cdot \text{sign}\!\left( \nabla_{\tilde{x}^{(t)}} \mathcal{L}_{\text{CE}}(f_\theta(\tilde{x}^{(t)}), y) \right) \right) \right).
\end{equation}

\subsubsection{PGD Adversarial Training (PGD-AT)}
To investigate whether adversarial training improves both PGD robustness and spatial explanation stability, we fine-tune ResNet-50 and Swin-Tiny using a mixed clean and 5-step PGD ($\epsilon=0.03$) min-max objective \cite{madry2018towards,zhang2019theoretically}:
\begin{equation}
    \min_\theta \mathbb{E}_{(\tilde{x}, y) \sim \mathcal{D}} \left[ \lambda \mathcal{L}_{\text{CE}}(f_\theta(\tilde{x}), y) + (1 - \lambda) \mathcal{L}_{\text{CE}}(f_\theta(\tilde{x}_{\text{PGD}}), y) \right],
\end{equation}
with $\lambda = 0.5$.

\subsection{Statistical Significance Testing}
To determine whether accuracy differences between any two classifiers $A$ and $B$ on the same test set are statistically significant, we apply McNemar's paired test \cite{mcnemar1947note} with continuity correction:
\begin{equation}
    \chi^2 = \frac{(|n_{10} - n_{01}| - 1)^2}{n_{10} + n_{01}},
\end{equation}
where $n_{10}$ is the number of test instances classified correctly by $A$ but misclassified by $B$, and $n_{01}$ is the converse, alongside the exact two-sided binomial $p$-value on discordant pairs.

\section{Experimental Setup}
\label{sec:setup}

\subsection{Datasets and Evaluation Protocol}
\begin{itemize}
    \item \textbf{Intel Natural Scenes:} Contains $\sim$17,000 natural scene images across $K=6$ categories (\emph{buildings}, \emph{forest}, \emph{glacier}, \emph{mountain}, \emph{sea}, \emph{street}). The official training split (14,034 images) is partitioned via class-stratified sampling into 85\% training (11,931 images) and 15\% validation (2,103 images) for each random seed $s \in \{0, 1, 2\}$, while the official held-out test set (3,000 images) is used exclusively for final evaluation.
    \item \textbf{SUN397-100 Large-Scale Protocol \cite{xiao2010sun}:} Because full 397-class multi-seed training paired with multi-method CAM and 10-step PGD evaluation exceeds single-session cloud GPU budgets, we define a deterministic, reproducible 100-class subset of SUN397 (selecting the top 100 categories by sample count, split 70\%/15\%/15\% into train/validation/test with exact class lists and indices stored in \texttt{results/config.json}). Classification is evaluated across the full held-out test split, while XAI faithfulness and FGSM/PGD robustness are evaluated on a fixed, class-balanced stratified subset of test images.
\end{itemize}

\subsection{Implementation Details and Hyperparameters}
All experiments are implemented in PyTorch 2.x using \texttt{timm}, \texttt{open\_clip}, and \texttt{pytorch-grad-cam}, and executed on NVIDIA Tesla T4 GPUs (16~GB VRAM). Table~\ref{tab:hyperparams} summarizes the training and evaluation hyperparameters.

\begin{table}[!t]
\centering
\caption{Summary of Experimental Hyperparameters.}
\label{tab:hyperparams}
\begin{tabular}{ll}
\hline
\textbf{Parameter} & \textbf{Value / Configuration} \\
\hline
Input Resolution & $224 \times 224$ RGB \\
Train Augmentations & RandomResizedCrop(0.75--1.0), Flip, ColorJitter(0.15) \\
Optimizer & AdamW ($\beta_1=0.9, \beta_2=0.999$, weight decay $10^{-4}$) \\
Learning Rate & $2 \times 10^{-4}$ (Backbones), $10^{-3}$ (CLIP Linear Probe) \\
LR Schedule & CosineAnnealingLR, Mixed Precision (AMP) \\
Epochs & 8 (Backbones), 5 (CLIP Linear), 4 (PGD-AT $\epsilon=0.03$) \\
Batch Size & 32 \\
Loss Function & Cross-Entropy with Label Smoothing ($\varepsilon_{\text{ls}}=0.05$) \\
Random Seeds & $\{0, 1, 2\}$ (Mean $\pm$ Std reported) \\
XAI Perturbation & 20 steps, $21\times21$ average-pooled blurred baseline \\
PGD-$L_\infty$ Attack & $T=10$ steps, $\alpha = 2.5\epsilon / T$, $\epsilon \in \{0.01, 0.03, 0.05\}$ \\
\hline
\end{tabular}
\end{table}

\section{Experimental Results}
\label{sec:results}

\subsection{Classification Accuracy and Computational Efficiency}
Table~\ref{tab:intel_prelim} reports our baseline single-seed classification performance and computational efficiency on Intel Natural Scenes, while Table~\ref{tab:intel_multiseed} presents the full 3-seed ($\{0, 1, 2\}$) comparison including Zero-Shot and Linear-Probe CLIP ViT-B/16. Swin-Tiny achieves the highest baseline accuracy of 94.60\% (Macro-F1: 94.69\%), outperforming ResNet-50 (93.97\%) and EfficientNet-B0 (93.87\%). Notably, EfficientNet-B0 matches ResNet-50 within 0.10\% accuracy while requiring $5.85\times$ fewer parameters (4.02M vs.\ 23.52M), whereas ResNet-50 achieves the lowest single-image inference latency (5.96~ms vs.\ 11.34~ms for Swin-Tiny).

\begin{table*}[!t]
\centering
\caption{Baseline Classification Performance, Efficiency, XAI Faithfulness, and PGD Robustness on Intel Natural Scenes (Canonical Run).}
\label{tab:intel_prelim}
\begin{tabular}{lcccccccc}
\hline
\textbf{Architecture} & \textbf{Test Acc. (\%)} & \textbf{Macro-F1 (\%)} & \textbf{Weighted-F1 (\%)} & \textbf{Params (M)} & \textbf{Latency (ms)} & \textbf{Grad-CAM (Ins$\uparrow$/Del$\downarrow$)} & \textbf{Grad-CAM++ (Ins$\uparrow$/Del$\downarrow$)} & \textbf{PGD ($\epsilon=0.03$)} \\
\hline
Swin-Tiny (ViT) & \textbf{94.60} & \textbf{94.69} & \textbf{94.57} & 27.52 & 11.34 & 0.6031 / 0.6710 & 0.7666 / \textbf{0.4899} & 1.67\% \\
ResNet-50 (CNN) & 93.97 & 94.11 & 93.94 & 23.52 & \textbf{5.96} & \textbf{0.8405} / 0.5538 & 0.8326 / 0.5664 & \textbf{19.17\%} \\
EfficientNet-B0 (CNN) & 93.87 & 94.00 & 93.85 & \textbf{4.02} & 8.24 & 0.7441 / 0.6560 & 0.6751 / 0.7441 & 0.00\% \\
\hline
\end{tabular}
\end{table*}

\begin{table*}[!t]
\centering
\caption{Multi-Seed ($s \in \{0, 1, 2\}$) Benchmark on Intel Natural Scenes Including CLIP ViT-B/16 Vision-Language Baselines (Mean $\pm$ Std).}
\label{tab:intel_multiseed}
\begin{tabular}{lcccccc}
\hline
\textbf{Architecture} & \textbf{Paradigm} & \textbf{Test Acc. (\%)} & \textbf{Macro-F1 (\%)} & \textbf{Grad-CAM Ins$\uparrow$ / Del$\downarrow$} & \textbf{Grad-CAM++ Ins$\uparrow$ / Del$\downarrow$} & \textbf{PGD ($\epsilon=0.03$) (\%)} \\
\hline
ResNet-50 & CNN & \textbf{[TBD]} & \textbf{[TBD]} & \textbf{[TBD]} & \textbf{[TBD]} & \textbf{[TBD]} \\
EfficientNet-B0 & CNN & \textbf{[TBD]} & \textbf{[TBD]} & \textbf{[TBD]} & \textbf{[TBD]} & \textbf{[TBD]} \\
Swin-Tiny & Hierarchical ViT & \textbf{[TBD]} & \textbf{[TBD]} & \textbf{[TBD]} & \textbf{[TBD]} & \textbf{[TBD]} \\
CLIP ViT-B/16 (Zero-Shot) & VLM (Frozen) & \textbf{[TBD]} & \textbf{[TBD]} & \textbf{[TBD]} & \textbf{[TBD]} & \textbf{[TBD]} \\
CLIP ViT-B/16 (Linear Probe) & VLM (Linear) & \textbf{[TBD]} & \textbf{[TBD]} & \textbf{[TBD]} & \textbf{[TBD]} & \textbf{[TBD]} \\
\hline
\end{tabular}
\end{table*}

\subsection{Quantitative XAI Faithfulness Analysis}
As shown in Table~\ref{tab:intel_prelim}, evaluating Insertion and Deletion AUC reveals a striking architectural divergence between convolutional and self-attention feature maps:
\begin{enumerate}
    \item \textbf{CNN Spatial Localization Advantage:} ResNet-50 achieves the highest overall Insertion AUC across all three attribution methods (0.8405 for Grad-CAM, 0.8326 for Grad-CAM++, and 0.8304 for LayerCAM), outperforming EfficientNet-B0 (0.7441) and Swin-Tiny (0.7666). Because ResNet-50 preserves continuous 2D spatial topology through residual bottlenecks, its final $7\times7$ feature maps align tightly with foreground scene structures.
    \item \textbf{Why Grad-CAM++ Outperforms Grad-CAM on Swin-Tiny:} On Swin-Tiny, standard Grad-CAM yields a comparatively low Insertion AUC of 0.6031 and Deletion AUC of 0.6710. However, switching to Grad-CAM++ boosts Insertion AUC by $+0.1635$ (to 0.7666) and lowers Deletion AUC by $-0.1811$ (to 0.4899---the best Deletion score across all models). In shifted-window self-attention, token activations in the final normalization layer distribute semantic information non-uniformly across window partitions; uniform global average pooling of gradients in standard Grad-CAM averages out localized positive token gradients, whereas the second-order weighting in Grad-CAM++ isolates dominant discriminative tokens.
\end{enumerate}

\subsection{Adversarial Robustness, Label-Leaking Diagnosis, and PGD-AT}
In our initial unstratified FGSM evaluation (Table~\ref{tab:intel_prelim}), single-step true-label FGSM exhibited non-monotonic accuracy on ResNet-50 (69.79\% at $\epsilon=0.01$ vs.\ 75.62\% at $\epsilon=0.03$) and EfficientNet-B0 (27.71\% vs.\ 70.63\%). Our code audit identified two compounding causes: (i) evaluating the first 15 unshuffled test batches sampled 91\% \emph{buildings} images (class 0), where sharp rectilinear edges amplify gradient sign correlation, and (ii) training with label smoothing ($\varepsilon_{\text{ls}}=0.05$, Eq.~\ref{eq:ls_ce}) induces the \emph{label leaking} effect \cite{kurakin2017adversarial}, where a large single-step true-label gradient step overshoots the narrow smoothed loss peak. Replacing true labels with predicted labels ($\text{FGSM}_{\text{pred}}$) or applying iterative 10-step $\text{PGD-}L_\infty$ completely eliminates label leaking, yielding strictly monotonic accuracy degradation across all models.

Under iterative PGD ($\epsilon=0.03$), a sharp robustness gap emerges: EfficientNet-B0 drops to 0.00\% and Swin-Tiny falls to 1.67\%, whereas ResNet-50 retains 19.17\% accuracy. Table~\ref{tab:pgd_at_results} and Table~\ref{tab:sun397_results} report the impact of PGD adversarial training (PGD-AT, $\epsilon=0.03$) and large-scale generalization on SUN397-100.

\begin{table}[!t]
\centering
\caption{Impact of PGD Adversarial Training ($\text{PGD-AT}, \epsilon=0.03$) on Clean Accuracy, PGD Robustness, and XAI Faithfulness (Intel Natural Scenes).}
\label{tab:pgd_at_results}
\begin{tabular}{lcccc}
\hline
\textbf{Model} & \textbf{Clean Acc.} & \textbf{PGD (.03)} & \textbf{Grad-CAM Ins$\uparrow$} & \textbf{Grad-CAM++ Ins$\uparrow$} \\
\hline
ResNet-50 (Clean) & 93.97\% & 19.17\% & 0.8405 & 0.8326 \\
ResNet-50 + PGD-AT & \textbf{[TBD]} & \textbf{[TBD]} & \textbf{[TBD]} & \textbf{[TBD]} \\
\hline
Swin-Tiny (Clean) & 94.60\% & 1.67\% & 0.6031 & 0.7666 \\
Swin-Tiny + PGD-AT & \textbf{[TBD]} & \textbf{[TBD]} & \textbf{[TBD]} & \textbf{[TBD]} \\
\hline
\end{tabular}
\end{table}

\begin{table*}[!t]
\centering
\caption{Large-Scale Benchmark Results on the 100-Class SUN397 Protocol (Classification, XAI Faithfulness, and Adversarial Robustness).}
\label{tab:sun397_results}
\begin{tabular}{lcccccc}
\hline
\textbf{Architecture} & \textbf{Test Acc. (\%)} & \textbf{Macro-F1 (\%)} & \textbf{Grad-CAM Ins$\uparrow$ / Del$\downarrow$} & \textbf{Grad-CAM++ Ins$\uparrow$ / Del$\downarrow$} & \textbf{FGSM (.03) (\%)} & \textbf{PGD (.03) (\%)} \\
\hline
ResNet-50 (CNN) & \textbf{[TBD]} & \textbf{[TBD]} & \textbf{[TBD]} & \textbf{[TBD]} & \textbf{[TBD]} & \textbf{[TBD]} \\
EfficientNet-B0 (CNN) & \textbf{[TBD]} & \textbf{[TBD]} & \textbf{[TBD]} & \textbf{[TBD]} & \textbf{[TBD]} & \textbf{[TBD]} \\
Swin-Tiny (ViT) & \textbf{[TBD]} & \textbf{[TBD]} & \textbf{[TBD]} & \textbf{[TBD]} & \textbf{[TBD]} & \textbf{[TBD]} \\
CLIP ViT-B/16 (Zero-Shot) & \textbf{[TBD]} & \textbf{[TBD]} & \textbf{[TBD]} & \textbf{[TBD]} & \textbf{[TBD]} & \textbf{[TBD]} \\
CLIP ViT-B/16 (Linear Probe) & \textbf{[TBD]} & \textbf{[TBD]} & \textbf{[TBD]} & \textbf{[TBD]} & \textbf{[TBD]} & \textbf{[TBD]} \\
\hline
\end{tabular}
\end{table*}

\section{Discussion}
\label{sec:discussion}

\subsection{The Accuracy--Faithfulness--Robustness Trilemma}
Our benchmark demonstrates that selecting a scene recognition backbone solely by top-1 clean accuracy masks critical operational trade-offs. While hierarchical Vision Transformers (Swin-Tiny) and multimodal foundation models (CLIP ViT-B/16) excel in clean accuracy due to long-range contextual aggregation and web-scale semantic pretraining, standard residual CNNs (ResNet-50) provide markedly higher spatial attribution faithfulness under first-order Grad-CAM and stronger native resilience against iterative PGD perturbations. Lightweight depthwise-separable architectures (EfficientNet-B0), although highly parameter-efficient (4.02M parameters), are exceptionally fragile under white-box PGD attacks (collapsing to 0.00\% at $\epsilon=0.01$), making them unsuitable for untrusted environments without adversarial hardening.

\subsection{Failure Mode Analysis and Practical Recommendations}
Qualitative inspection of misclassified test samples with Grad-CAM overlays reveals that residual errors cluster primarily along semantically overlapping natural boundaries---specifically between \emph{glacier} and \emph{mountain} (where snow-covered rocky peaks share identical visual textures) and between urban \emph{street} and \emph{buildings} scenes. Based on our findings, we recommend: (1) pairing Swin-Tiny or ViT backbones with \textbf{Grad-CAM++} or \textbf{LayerCAM} rather than vanilla Grad-CAM for visual auditing; (2) using \textbf{class-stratified subsets} and \textbf{iterative PGD / $\text{FGSM}_{\text{pred}}$} when auditing label-smoothed classifiers to avoid label-leaking artifacts; and (3) applying mixed-objective \textbf{PGD-AT} when deploying scene classifiers in security-sensitive pipelines.

\subsection{Limitations}
To operate within single-session cloud GPU constraints (2$\times$ NVIDIA T4), our large-scale SUN397 evaluation adopts a deterministic 100-class subset protocol and evaluates pixel-perturbation AUC and PGD curves on a fixed, class-balanced test subset of 500--1,000 images. Additionally, our CLIP experiments focus on zero-shot prompt ensembling and frozen-backbone linear probing rather than end-to-end multimodal contrastive fine-tuning.

\section{Conclusion and Future Work}
\label{sec:conclusion}
This paper presented a unified benchmark of CNNs (ResNet-50, EfficientNet-B0), a hierarchical Vision Transformer (Swin-Tiny), and a Vision-Language Model (CLIP ViT-B/16) for explainable multi-class scene recognition. By jointly measuring predictive performance, Insertion/Deletion XAI faithfulness, white-box FGSM/PGD robustness, and the effect of PGD adversarial training across Intel Natural Scenes and SUN397-100, we exposed a clear trade-off between clean accuracy, spatial explanation quality, and adversarial resilience.

In future work, we plan to extend this explainable robustness framework from holistic scene categorization to **natural language-driven, open-vocabulary object detection and visual grounding** using architectures such as Grounding DINO \cite{liu2024grounding}, enabling fine-grained, text-conditioned spatial attribution and adversarial defense at both the object and scene levels.

\section*{Data and Code Availability}
All source code, data split configurations, automated LaTeX table scripts, and interactive Gradio demo files are openly available at \url{https://github.com/matiur-rahman1/explainable-scene-recognition}.

\bibliographystyle{IEEEtran}
\bibliography{references}

\end{document}
